
\documentclass[12pt,a4paper]{article}

\usepackage{amsmath}
\usepackage{amssymb}
\usepackage{amsthm}
\usepackage{natbib}

\newtheorem{theorem}{Theorem}[section]

\begin{document}
	
	\title{My First LaTeX Document}
	\author{Your Name}
	\date{\today}
	
	\maketitle
	
	\tableofcontents
	
	\newpage
	
	\section{Introduction}
	
	This is my first LaTeX document.
	
	Words are separated by one or more spaces. Paragraphs are separated by one or more blank lines.
	
	\section{Background}
	
	In Section \ref{intro}, we discussed the introduction.
	
	Double quotes are typed like this: ``quoted text''.
	
	Single quotes are typed like this: `single-quoted text'.
	
	\textbf{Long dashes} are typed as three dash characters---like this.
	
	Italic text is typed like this: \textit{This is italic text}.
	
	Bold text is typed like this: \textbf{This is bold text}.
	
	\subsection{Mathematical Expressions}
	
	{\large $\beta = a^2 + b\times c + 2$}
	
	An example of a mathematical expression is $a^2+b^2=c^2$.
	
	\subsection{Numbered List}
	
	\begin{enumerate}
		\item First item
		\item Second item
		\item Third item
	\end{enumerate}
	
	\subsection{Theorem}
	
	\begin{theorem}
		There are infinitely many prime numbers.
	\end{theorem}
	
	As stated in Theorem \ref{T1}, there are infinitely many prime numbers.
	
	\section{A Warning or Two}
	
	If you get too much space after a mid-sentence period---abbreviations like etc.\ are the common culprits---then type a backslash followed by a space after the period.
	
	\noindent Remember, special characters such as dollar sign and backslash must be used carefully in LaTeX.
	
	The following characters are printed by typing a backslash in front of them:
	
	\$ \quad \& \quad \# \quad \% \quad \_ \quad \{ \quad \}
	
	\section{Science!}
	\label{intro}
	
	Einstein's famous mass-energy equivalence equation is:
	
	$E=mc^2$
	
	This equation describes the relationship between mass and energy.
	
	Einstein introduced this relationship in his 1905 research paper \citep{einstein1905}.
	
	\subsection{Schrödinger Equation}
	
	The time-dependent Schrödinger equation is a fundamental equation of quantum mechanics.
	
	It is given by:
	
	\begin{equation}
		i\hbar\frac{\partial\Psi(\mathbf{r},t)}{\partial t}
		=
		\left[
		-\frac{\hbar^2}{2\mu}\nabla^2
		+V(\mathbf{r},t)
		\right]\Psi(\mathbf{r},t)
		\label{eq:schrodinger}
	\end{equation}
	
	As shown in Equation \ref{eq:schrodinger}, the Schrödinger equation describes the evolution of the quantum wave function over time.
	
	Schrödinger introduced this equation in his 1926 research paper \citep{schrodinger1926}.
	
	\newpage
	
	\bibliographystyle{agu}
	\bibliography{physics}
	
\end{document}
